% Lösungen Winkel und Dreiecke (zusammenbau v0.4, Heft)
\einheitenkopf{Winkel und Dreiecke}

\erg{48}{a) spitz \quad b) $\alpha = 37^\circ$ \quad c) $\alpha = 118^\circ$}
\erg{49}{a) $\gamma = 44^\circ - 41^\circ = 3^\circ$ \quad b) $\gamma = 29^\circ - 22^\circ = 7^\circ$ \quad c) 730 m + 684 m = 1\,414 m \quad d) 615 m + 470 m = 1\,085 m}
\erg{50}{a) Scheitelwinkel \quad b) $38^\circ$ (Scheitelwinkel sind gleich groß) \quad c) $180^\circ - 44^\circ = 136^\circ$}
\erg{51}{a) Scheitelwinkel $64^\circ$; $\beta = 64^\circ - 23^\circ = 41^\circ$ \quad b) $49^\circ$ (Wechselwinkel) \quad c) $\beta = 43^\circ$ (Scheitelwinkel); die Angabe $58^\circ$ wird nicht gebraucht \quad d) $\alpha = 180^\circ - 107^\circ = 73^\circ$ (oder $360^\circ - 58^\circ - 122^\circ - 107^\circ$) \quad e) $\gamma = 180^\circ - 52^\circ = 128^\circ$}
\erg{52}{a) Stufenwinkel an g $44^\circ$; $\alpha = 180^\circ - 44^\circ = 136^\circ$ \quad b) Stufenwinkel an g $51^\circ$; $\alpha = 180^\circ - 51^\circ = 129^\circ$ \quad c) $\beta = 180^\circ - 79^\circ = 101^\circ$ \quad d) $\beta = 180^\circ - 83^\circ = 97^\circ$}
\erg{53}{a) rechtwinklig \quad b) $\gamma = 180^\circ - 52^\circ - 71^\circ = 57^\circ$ \quad c) $180^\circ - 38{,}6^\circ - 94{,}7^\circ = 46{,}7^\circ$}
\erg{54}{a) $180^\circ - 51{,}4^\circ - 76{,}9^\circ = 51{,}7^\circ$ \quad b) im Teildreieck ADC: $\gamma_1 = 180^\circ - 90^\circ - 27{,}6^\circ = 62{,}4^\circ$ \quad c) $\angle ACB = 180^\circ - 90^\circ - 22^\circ = 68^\circ$; Nebenwinkel $\angle ACD = 180^\circ - 68^\circ = 112^\circ$; $\delta = 180^\circ - 112^\circ - 9^\circ = 59^\circ$ \quad d) $\delta = 180^\circ - 114^\circ - 38^\circ = 28^\circ$ \quad e) Im Dreieck ABF: rechter Winkel bei F, $45^\circ$ bei A, also bei B $180^\circ - 90^\circ - 45^\circ = 45^\circ$; gleiche Winkel bei A und B, also ist ABF gleichschenklig mit BF = AF. \quad f) gibt es zwei Paare gleich langer Nachbarseiten \quad g) sind gegenüberliegende Winkel gleich groß}
\erg{55}{a) BC = AC = 7 cm (gleiche Winkel bei A und B); $\gamma = 180^\circ - 2 \cdot 66^\circ = 48^\circ$ \quad b) BC = AC = 2,4 m (gleiche Winkel bei A und B); $\gamma = 180^\circ - 2 \cdot 73^\circ = 34^\circ$ \quad c) AF = FC = 18 cm = DF. Die Dreiecke AFD und CFD sind gleichschenklig und rechtwinklig bei F, ihre Basiswinkel messen je $45^\circ$; bei D also $45^\circ + 45^\circ = 90^\circ$. (Auch: D liegt auf dem Kreis um F mit Radius 18 cm, AC ist sein Durchmesser – Thales.) \quad d) AF = FC = 22 cm = DF. Die Dreiecke AFD und CFD sind gleichschenklig und rechtwinklig bei F, ihre Basiswinkel messen je $45^\circ$; bei D also $45^\circ + 45^\circ = 90^\circ$. (Auch Thales: D liegt auf dem Kreis mit dem Durchmesser AC.)}
\erg{56}{a) SWS \quad b) c = 7 cm von A nach B; Kreisbogen um A mit 5 cm, Kreisbogen um B mit 6 cm; der Schnittpunkt ist C. Kontrolle: $\gamma \approx 78^\circ$ \quad c) c = 6 cm von A nach B; in A den Winkel $50^\circ$ antragen; auf dem freien Schenkel b = 5 cm abtragen, Endpunkt C. Kontrolle: a ≈ 4,7 cm}

56b)

\dreieck{(0,0)}{(7,0)}{(2.71,4.20)}{a}{b}{c}{\alpha}{\beta}{\gamma}


56c)

\dreieck{(0,0)}{(6,0)}{(3.21,3.83)}{a}{b}{c}{\alpha}{\beta}{\gamma}

\erg{57}{a) $b + c > 12$ m \quad b) $b + c > 15$ cm}
